\documentclass[11pt]{article}
\usepackage[T1]{fontenc}
\usepackage{amsmath,amssymb,amsthm}
\usepackage[margin=1in]{geometry}
\usepackage[hidelinks]{hyperref}
\newtheorem{theorem}{Theorem}
\newtheorem{corollary}[theorem]{Corollary}
\newtheorem{lemma}[theorem]{Lemma}
\newtheorem{proposition}[theorem]{Proposition}
\theoremstyle{definition}
\newtheorem{remark}[theorem]{Remark}
\DeclareMathOperator{\ad}{ad}
\DeclareMathOperator{\im}{im}
\DeclareMathOperator{\End}{End}
\DeclareMathOperator{\Hom}{Hom}
\DeclareMathOperator{\Tri}{Tri}
\title{Unique addition in centerless metabelian Lie algebras of finite length}
\author{}
\date{September 8, 2026}
\begin{document}
\maketitle
\begin{abstract}
We prove that a centerless metabelian Lie algebra over a commutative unital ring is a unique-addition Lie ring whenever its underlying module has finite length. Thus every bracket-preserving bijection from it onto an arbitrary Lie ring is additive. The theorem includes all finite-dimensional centerless metabelian Lie algebras over arbitrary fields and all finite centerless metabelian Lie rings. The proof first recovers addition on the abelian Fitting images of individual adjoint operators, then eliminates the remaining defect using commuting operators on the centralizer of the derived algebra. Over each finite field we exhibit an algebra covered by the theorem for which every adjoint action on the derived algebra is singular. We also give families with unbounded common-centralizer number and a separate faithful-action criterion without a finiteness assumption.
\end{abstract}

\section{Introduction}
A Lie ring is an abelian group with a biadditive alternating bracket satisfying the Jacobi identity. A Lie ring $L$ has \emph{unique addition}, or is a \emph{UA-Lie ring}, if every bijection $\alpha:L\to S$ onto an arbitrary Lie ring that satisfies
\[
 \alpha([x,y])=[\alpha(x),\alpha(y)]\qquad(x,y\in L)
\]
is additive. When $L$ is an algebra over a ring $R$, no $R$-module structure on $S$ is required. We use the alternating convention for Lie algebras over rings, including in characteristic two.

Arzhantsev~\cite{Arz} studies this reconstruction property and asks whether a centerless Lie ring can fail to have unique addition. His two-centralizer criterion covers many finite-dimensional algebras over infinite fields. Example~2 and Problem~4 concern a centerless solvable matrix algebra outside that criterion. Our main result treats an entire class over arbitrary fields and, more generally, over rings.

\begin{theorem}\label{thm:main}
Let $R$ be a commutative unital ring and let $L$ be a Lie algebra over $R$. Suppose that the underlying $R$-module of $L$ has finite length, that $Z(L)=0$, and that $[L,L]$ is abelian. Then $L$ is a UA-Lie ring.
\end{theorem}

Here $[L,L]$ denotes the $R$-submodule generated by all brackets; it also equals the derived Lie subring, since $r[x,y]=[rx,y]$. An algebra with abelian derived algebra is called \emph{metabelian}.

\begin{corollary}\label{cor:fields}
Every finite-dimensional centerless metabelian Lie algebra over an arbitrary field is a UA-Lie ring. Every finite centerless metabelian Lie ring is a UA-Lie ring.
\end{corollary}
\begin{proof}
For the first assertion take $R$ to be the field. For the second take $R=\mathbb Z$; every finite abelian group has finite length as a $\mathbb Z$-module.
\end{proof}

There is important precedent for special cases. Qi--Hou~\cite{QH} proved almost additivity for Lie-multiplicative bijections on triangular associative algebras; Ferreira--Guzzo~\cite{FG} reproduce and extend their theorem. For a unital associative algebra $A$ acting faithfully on $V$, the associated affine Lie algebra is a central factor of $\Tri(A,V,K)$. The earlier proof adapts to arbitrary Lie-ring targets and gives unique addition for this affine case, including the particular matrix algebra in Arzhantsev's question. We explain this overlap in Section~\ref{sec:comparison}; the particular matrix answer is not our principal novelty claim.

Theorem~\ref{thm:main} requires neither an associative acting algebra nor a single element acting invertibly on the derived algebra. The latter distinction matters over finite fields: Section~\ref{sec:examples} gives examples where all such actions are singular. Our proof uses several intrinsic Fitting images instead of one global splitting. Sosnov's recent classification for $\mathfrak{gl}_n(K)$ and $\mathfrak{sl}_n(K)$~\cite{Sos} concerns different classes of Lie rings.

\section{Addition on an abelian Fitting image}
All subtractions involving $\alpha$ below occur in the given additive group of its target. First observe that every bracket-preserving map sends zero to zero, since
\[
 \alpha(0)=\alpha([0,0])=[\alpha(0),\alpha(0)]=0.
\]
For a bracket-preserving bijection, centers correspond as sets.

\begin{lemma}\label{lem:fitting}
Let $L$ be a Lie ring, let $e\in L$, and put $d=\ad_e$. Suppose that some integer $N\ge1$ satisfies
\[
 \im d^{N+1}=\im d^N,\qquad \ker d^{2N}=\ker d^N.
\]
If $U=\im d^N$ is abelian, then every bracket-preserving bijection $\alpha:L\to S$ onto a Lie ring is additive on $U$.
\end{lemma}
\begin{proof}
Put $W=\ker d^N$. The first hypothesis gives $d(U)=U$. If $u=d^Nx$ belongs to $U\cap W$, then $d^{2N}x=0$, so $x\in\ker d^N$ and $u=0$. Therefore $U\cap W=0$, and $d$ restricts bijectively to $U$.

The subgroup $W$ is closed under brackets. Indeed, if $a,b\in W$, the iterated derivation identity gives
\[
 d^{2N}[a,b]=\sum_{j=0}^{2N}\binom{2N}{j}[d^ja,d^{2N-j}b]=0.
\]
In every term at least one exponent is at least $N$. The kernel hypothesis implies $[a,b]\in W$.

Set $E=\alpha(e)$ and $D=\ad_E$. For every $j\ge0$ we have $D^j\alpha(x)=\alpha(d^jx)$. Consequently
\[
 U'=\alpha(U)=\im D^N,\qquad W'=\alpha(W)=\ker D^N.
\]
Both are additive subgroups of $S$. The subgroup $U'$ is abelian, $W'$ is bracket closed, $U'\cap W'=0$, and $D$ restricts bijectively to $U'$. These conclusions follow from the source properties and set conjugacy; additivity of $\alpha$ is not being assumed.

For $v\in U$, define
\[
 B_v=\alpha(e+v)-\alpha(v).
\]
Since $N\ge1$ and $d(e)=0$, we have $d^N(e+v)=d^Nv$, so $B_v\in W'$. For $w\in U$, abelianness of $U$ gives
\begin{align*}
 [B_v,\alpha(w)]
 &=[\alpha(e+v),\alpha(w)]-[\alpha(v),\alpha(w)]\\
 &=\alpha([e,w])=D\alpha(w).
\end{align*}
Now $[e+v,e+w]=d(w-v)$. Expanding the corresponding target bracket yields
\[
 \alpha(d(w-v))=[B_v,B_w]+D(\alpha(w)-\alpha(v)).
\]
The left side and the last term lie in $U'$, whereas $[B_v,B_w]\in W'$. Thus the middle term vanishes, and
\[
 D\alpha(w-v)=D(\alpha(w)-\alpha(v)).
\]
Injectivity of $D$ on $U'$ gives $\alpha(w-v)=\alpha(w)-\alpha(v)$. This implies additivity on $U$.
\end{proof}

\begin{lemma}\label{lem:nilpotents}
Let $T_1,\ldots,T_m$ be commuting endomorphisms of an abelian group $C$ with $\bigcap_i\ker T_i=0$. For every integer $N\ge1$,
\[
 \bigcap_i\ker T_i^N=0.
\]
\end{lemma}
\begin{proof}
Let $Q=\bigcap_i\ker T_i^N$. Commutativity makes $Q$ invariant under every $T_i$, and each $T_i$ is nilpotent on $Q$. If $Q\ne0$, successively restrict to the kernels of $T_1,\ldots,T_m$. At each step the current subgroup is nonzero and invariant under the remaining operators. A nilpotent endomorphism of a nonzero group cannot be injective, so its kernel is nonzero. This produces a nonzero element of $\bigcap_i\ker T_i$, a contradiction.
\end{proof}

\section{Proof of the finite-length theorem}
\begin{proof}[Proof of Theorem~\ref{thm:main}]
The zero algebra is immediate, so suppose $L\ne0$. Let $\ell$ be its $R$-module length and fix $N\ge\max(1,\ell)$. For every $R$-linear endomorphism $d$ of $L$, the chains of kernels and images stabilize by exponent $N$: each strict inclusion changes length, and one equality forces all later terms to agree. In particular,
\[
 \im d^{N+1}=\im d^N,\qquad \ker d^{2N}=\ker d^N.
\]

Put $V=[L,L]$ and $C=C_L(V)$. The centralizer $C$ is an $R$-submodule and an ideal. To see invariance, let $x\in L$, $c\in C$, and $v\in V$. The Jacobi identity and $[x,v]\in V$ give $[[x,c],v]=0$.

A finite-length module is finitely generated, by induction on a composition series whose simple factors are cyclic. Choose $R$-module generators $x_1,\ldots,x_m$ for $L$. The restrictions
\[
 T_i=\ad_{x_i}|_C
\]
commute: $[T_i,T_j]=\ad_{[x_i,x_j]}|_C=0$, since $[x_i,x_j]\in V$. Their common kernel is zero. Indeed, an element of $C$ commuting with every generator commutes, by $R$-bilinearity, with all of $L$ and hence is central. Lemma~\ref{lem:nilpotents} gives
\begin{equation}\label{eq:joint}
 C\cap\bigcap_{i=1}^m\ker(\ad_{x_i})^N=0.
\end{equation}

Fix a bracket-preserving bijection $\alpha:L\to S$. For every $e\in L$, its Fitting image
\[
 U_e=\im(\ad_e)^N
\]
lies in $V$ and is therefore abelian. Lemma~\ref{lem:fitting} shows that $\alpha$ is additive on each $U_e$.

We first recover addition on $V$. For $v,w\in V$, form the defect
\[
 \delta=\alpha(v+w)-\alpha(v)-\alpha(w)\in S.
\]
Every element of $V$ commutes with $v,w,v+w$, so $\delta$ centralizes $\alpha(V)$. By surjectivity, $\delta=\alpha(a)$ for some $a\in L$. Bracket preservation and injectivity give $a\in C$.

Put $D_i=\ad_{\alpha(x_i)}$. The three source arguments in the next calculation belong to $U_{x_i}$, on which $\alpha$ is additive. Hence
\begin{align*}
 D_i^N\delta
 &=\alpha((\ad_{x_i})^N(v+w))\\
 &\quad-\alpha((\ad_{x_i})^Nv)-\alpha((\ad_{x_i})^Nw)=0.
\end{align*}
Intertwining implies $(\ad_{x_i})^Na=0$ for every $i$. Equation~\eqref{eq:joint} gives $a=0$, so $\delta=0$. Thus $\alpha$ is additive on $V$.

For arbitrary $x,y,z\in L$, all of $[x+y,z]$, $[x,z]$ and $[y,z]$ lie in $V$. Its recovered addition therefore gives
\[
 [\alpha(x+y)-\alpha(x)-\alpha(y),\alpha(z)]=0.
\]
Since $\alpha$ is surjective and $Z(S)=\alpha(Z(L))=0$, the defect vanishes. Thus $\alpha$ is additive on $L$.
\end{proof}

\begin{remark}
The commuting operators in the proof act on $C_L([L,L])$, not necessarily on all of $L$. No equality $C_L([L,L])=[L,L]$ is needed. Finite length provides stabilization and finitely many generators; neither property is asserted for arbitrary infinite-dimensional Lie algebras or arbitrary finitely generated modules over rings.
\end{remark}

\section{Examples and consequences}\label{sec:examples}
\subsection{Finite fields without a regular adjoint}
Let $K=\mathbb F_q$ and $A=K^2$. Choose a set $\Lambda$ containing one nonzero linear form from each one-dimensional subspace of $A^*$. Define
\[
 V=\bigoplus_{\lambda\in\Lambda}K v_\lambda,
 \qquad a(v_\lambda)=\lambda(a)v_\lambda,
 \qquad L_q=A\ltimes V.
\]
The algebra $A$ is abelian, and the displayed operators commute.

\begin{proposition}\label{prop:finite}
The algebra $L_q$ is a centerless metabelian UA-Lie ring of dimension $q+3$. Every element of $L_q$ acts singularly on its derived algebra.
\end{proposition}
\begin{proof}
The forms in $\Lambda$ span $A^*$, so the action is faithful. Their weight spaces have nonzero weights, so $V$ has no common fixed vector. These facts give centerlessness of $A\ltimes V$. Each $v_\lambda$ belongs to $[A,V]$, because $\lambda$ takes a nonzero value; hence $[L_q,L_q]=V$. There are $q+1$ forms, giving the dimension and, by Corollary~\ref{cor:fields}, the UA property.

For nonzero $a\in A$, its annihilator in $A^*$ is one-dimensional, so some $\lambda\in\Lambda$ vanishes on $a$. Thus $a$ acts singularly on $V$; this is also true for $a=0$. Adding a translation from $V$ does not change that action.
\end{proof}

Thus a proof requiring one invertible adjoint on the derived algebra cannot cover the full theorem. Over $\mathbb F_2$, the example has dimension five. No claim that $L_q$ fails the two-centralizer condition is needed here.

\subsection{Unbounded common-centralizer number}
For a finite-dimensional centerless Lie algebra $L$, let $c(L)$ be the least number of elements whose common centralizer is zero. Let $K$ be any field, $U=K^r$, $W=K^s$, and $V=U\oplus W$, where $r,s\ge1$. Extend each map in $\Hom_K(W,U)$ by zero on $U$, and set
\[
 A=K I_V\oplus\Hom_K(W,U),\qquad L_{r,s}=A\ltimes V.
\]
The acting algebra is abelian, since the product of any two of its off-diagonal operators is zero.

\begin{proposition}\label{prop:centralizers}
The algebra $L_{r,s}$ is centerless and metabelian, is a UA-Lie ring, and satisfies $c(L_{r,s})=s+1$.
\end{proposition}
\begin{proof}
Write arbitrary elements as $x_i=(a_i,M_i,u_i,w_i)$, with $a_i\in K$. An element $z=(0,N,v,0)$ centralizes $x_1,\ldots,x_h$ exactly when
\[
 N(w_i)-a_iv=0\qquad(1\le i\le h).
\]
Identify $(N,v)$ with the map $F:W\oplus K\to U$ defined by $F(w,a)=N(w)-av$. The solution space has dimension
\[
 r\bigl(s+1-\dim\langle(w_i,a_i):1\le i\le h\rangle\bigr)
 \ge r(s+1-h).
\]
It is nonzero if $h\le s$.

Conversely, choose a basis $w_1,\ldots,w_s$ of $W$. The common centralizer of $e=(I_V,0)$ and the translation vectors $w_1,\ldots,w_s$ is zero. Commuting with $e$ removes the translation component. If $aI_V+M$ annihilates every $w_j$, its $W$-components give $a=0$, and then $M=0$. This proves centerlessness and the upper bound. The algebra is metabelian by construction, so Corollary~\ref{cor:fields} gives unique addition.
\end{proof}

In block-matrix notation this family is
\[
 L_{r,s}=\left\{
 \begin{pmatrix}
 aI_r&M&u\\0&aI_s&w\\0&0&0
 \end{pmatrix}:a\in K,\ M\in\Hom_K(W,U),\ u\in U,\ w\in W
 \right\}.
\]
Taking $r=s=2$ gives Arzhantsev's Example~2. As discussed below, its UA property is also accessible by the older triangular method. The exact common-centralizer count quantifies the failure of a uniform bound on $c(L)$ among UA-Lie algebras.

Over $\mathbb F_q$, the direct sum $L_q\oplus L_{r,s}$ is again centerless and metabelian. Every adjoint on its derived algebra is singular by the first factor, while any $s$ elements have a nonzero common centralizer by the second. Theorem~\ref{thm:main} applies to these examples with both features simultaneously.

\subsection{Reconstruction consequences}
For any $L$ in Theorem~\ref{thm:main}, the bracket-preserving self-bijections are exactly its Lie-ring automorphisms. A bracket-preserving bijection to another Lie ring is a Lie-ring isomorphism. If source and target are Lie algebras over the same prime field, additivity also forces linearity. Over a nonprime field, the theorem does not assert scalar linearity or semilinearity, and it does not give a complexity bound for reconstructing addition from a bracket table.

\section{A separate criterion without finite length}
The following result applies to arbitrary abelian translation groups and allows a nonmetabelian acting Lie ring. It is independent in scope from Theorem~\ref{thm:main}.

\begin{theorem}\label{thm:affine}
Let $V$ be an abelian group and $A$ a Lie subring of $\End_{\mathbb Z}(V)$ with its commutator bracket. Suppose $A$ contains an element $t$ that commutes with every element of $A$ and is bijective on $V$. Then the Lie ring
\[
 L=A\ltimes V,\qquad
 [(a,v),(b,w)]=([a,b],a(w)-b(v)),
\]
has unique addition.
\end{theorem}
\begin{proof}
Let $\alpha:L\to S$ be a bracket-preserving bijection. Put $e=(t,0)$, $E=\alpha(e)$ and $D=\ad_E$. Identifying source summands with $A$ and $V$, we have
\[
 \im D=\alpha(V)=V',\qquad \ker D=\alpha(A)=A'.
\]
These are additive subgroups. The group $V'$ is abelian and $D$ restricts bijectively to it.

For $a\in A$ and $v\in V$, the identity $[e,a+v]=[e,v]$ implies
\[
 \alpha(a+v)-\alpha(v)=\alpha(b)\quad\text{for some }b\in A.
\]
Bracketing with $\alpha(w)$ for $w\in V$ gives $\alpha(b(w))=\alpha(a(w))$. Injectivity and faithfulness of the action imply $b=a$. Hence
\begin{equation}\label{eq:mixed}
 \alpha(a+v)=\alpha(a)+\alpha(v).
\end{equation}
Apply this to $[e+v,e+w]=t(w-v)$. It gives
\[
 D\alpha(w-v)=D(\alpha(w)-\alpha(v)).
\]
Injectivity on $V'$ proves additivity on $V$. For $a,b\in A$, write $\alpha(a)+\alpha(b)=\alpha(c)$ with $c\in A$. Bracketing with every $\alpha(w)$ and using additivity on $V$ yields $c(w)=(a+b)(w)$. Faithfulness gives $c=a+b$. Thus $\alpha$ is additive on $A$, and \eqref{eq:mixed} completes the proof.
\end{proof}

The identity endomorphism is an admissible choice of $t$. The acting Lie subring need not be closed under associative multiplication, and $t^{-1}$ need not belong to it. Theorem~\ref{thm:affine} imposes a global invertible action; Proposition~\ref{prop:finite} shows why Theorem~\ref{thm:main} requires a different argument.

\section{Prior methods and boundaries}\label{sec:comparison}
Suppose $A$ is an associative unital $K$-algebra acting faithfully on $V$. The underlying Lie ring of $\Tri(A,V,K)$ is isomorphic to
\[
 (A\ltimes V)\oplus K,
\]
where the last factor is central: send $(a,v,b)$ to $((a-bI_V,v),b)$. Extend a bracket-preserving bijection $\alpha:A\ltimes V\to S$ by the identity on $K$. The $n=2$ proof in Ferreira--Guzzo~\cite[Theorem 4.1]{FG}, extending Qi--Hou~\cite{QH}, uses only addition and brackets in the target. It therefore gives additivity modulo the center for this extension to the Lie-ring target $S\oplus K$. The source $A\ltimes V$ is centerless: commuting with the identity action kills its translation component, and faithfulness kills the remaining component. Thus $S$ is centerless as well. The additive defect lies in the extra central factor, but its extra coordinate is zero, so it vanishes.

This recovers the associative-unital affine case by adapting prior proofs. The finite-length theorem instead uses abelian Fitting images and their joint detection of defects, without associative closure or a global invertible action. It does not use inheritance of unique addition under arbitrary Lie subalgebras, quotients, or scalar extension.

The general centerless-Lie-ring problem remains beyond these results. For a nonmetabelian algebra, Fitting images need not be abelian. In infinite-dimensional situations, the relevant kernel and image chains need not stabilize. These are specific obstructions to the proof, not counterexamples to unique addition. Surjectivity of the bracket-preserving map is also essential to the reconstruction arguments given here.

\begin{thebibliography}{9}
\raggedright
\bibitem{Arz} I. Arzhantsev, \emph{Uniqueness of addition in Lie algebras revisited}, Math. Commun. \textbf{30} (2025), 179--189, \href{https://doi.org/10.64785/mc.30.2.3}{doi:10.64785/mc.30.2.3}.
\bibitem{FG} B. L. M. Ferreira and H. Guzzo Jr., \emph{Lie $n$-multiplicative mappings on triangular $n$-matrix rings}, Rev. Un. Mat. Argentina \textbf{60} (2019), no. 1, 9--20, \href{https://doi.org/10.33044/revuma.v60n1a02}{doi:10.33044/revuma.v60n1a02}.
\bibitem{QH} X. Qi and J. Hou, \emph{Additivity of Lie multiplicative maps on triangular algebras}, Linear Multilinear Algebra \textbf{59} (2011), no. 4, 391--397, \href{https://doi.org/10.1080/03081080903582094}{doi:10.1080/03081080903582094}.
\bibitem{Sos} G. Sosnov, \emph{Uniqueness of addition in Lie rings $\mathfrak{gl}_n(K)$ and $\mathfrak{sl}_n(K)$}, preprint (2026), \href{https://arxiv.org/abs/2606.08201}{arXiv:2606.08201v1}.
\end{thebibliography}
\end{document}
